\section{主張境界と証拠ゲート}

本稿では、外部文献、公開データセット、tool output、実験結果をすべて
evidence として扱い、単独では authority とみなさない。表
\ref{tab:claim-gates} に、今後の論文化で各主張を昇格するための最低条件を示す。
この表を満たさない文は、本文では仮説、計画、または評価項目として記述する。

\begin{table}[tb]
  \centering
  \caption{主張クラスと昇格に必要な証拠}
  \label{tab:claim-gates}
  \small
  \begin{tabularx}{\linewidth}{lYY}
    \toprule
    主張クラス & 必要な証拠 & 現在の扱い \\
    \midrule
    文献・既存研究 claim &
    source URL、DOI、可能なら source paper または publication page の監査記録。 &
    background evidence。novelty や SOTA は非主張。 \\
    dataset claim &
    schema、license、run 単位、label 意味、trace split 可能性の source audit。 &
    RanSAP/RanSMAP は候補。利用可能性は未確定。 \\
    性能 claim &
    experiment manifest、固定 split、feature schema、seed、metrics、verdict。 &
    非主張。評価プロトコルのみ。 \\
    false-positive claim &
    benign workload holdout、false alarms per volume per day、threshold tuning split。 &
    非主張。主要評価項目。 \\
    未知脅威 claim &
    AE は benign のみで学習・calibration し、攻撃 family は評価専用に保持する。
    教師あり参照は既知 family で学習し、held-out family で別集計する。 &
    非主張。family holdout は proxy 評価。 \\
    device-fit claim &
    MNN 変換済み model、offline/MNN score parity、volume あたりの model weight と
    入力統計/state、inference slot あたりの transient scratch、slot 数 $Q$、
    共有 weight と per-volume state の分解、volume 数 $V$ に対する aggregate memory と
    10 秒あたり CPU 時間。 &
    非主張。500 KB は per-volume detector-data 制約条件。 \\
    \bottomrule
  \end{tabularx}
\end{table}
